\scanpage{27}
\mspage{\handwritten{19}}
% Source viewed after a 180-degree rotation. Archive number has the opposite orientation.
% The whole draft is cancelled by the long source paths drawn over the complete reading below.
% Local overtyping/deletion remains independently encoded; the global cancellation is not a horizontal strike.
% Source coordinates in the graphic are those of the upright, 1685 x 2159 view.
% The marginal oval is kept inside this coordinate system to preserve its leftward and vertical placement.
\begingroup
\makeatletter\tagsleft@true\makeatother
\figure{\begin{tikzpicture}[x=.0072cm,y=-.0072cm,every node/.style={inner sep=0pt,font=\ttfamily\fontsize{8}{14.3}\selectfont}]
\path[use as bounding box] (55,180) rectangle (1650,2145);
\node[anchor=north west,text width=10.18cm,align=justify] at (230,200) {2.1.1.\hspace{.54cm}On peut dire encore qu'une structure de biextension sur le $G_{P\times Q}$-torseur $E$ est définie par deux lois de composition non partout définies $\handwritten{\phi}(x,y)$ et $\handwritten{\psi}(x,y)$ sur $E$, la première étant définie lorsque \deleted{$j($} $\mathrm{pr}_2j(x)=\mathrm{pr}_2j(y)$, la deuxième lorsque $\mathrm{pr}_1j(x)=\mathrm{pr}_1j(y)$ \deleted{(où} (où $j:E\handwritten{\longrightarrow}P\times Q$ est le morphisme structural du torseur). Les axiomes, en plus des associativités du type};

\node at (109,373) {\handwritten{$\unsure{\psi|\Psi}$}};
\draw[line width=.3pt] (118,325) .. controls (83,328) and (62,378) .. (71,411) .. controls (76,455) and (122,424) .. (138,378) .. controls (146,353) and (137,337) .. (130,338);

\node[anchor=north west,text width=6.2cm] at (455,612) {\[
\deleted{\handwritten{\phi}(gx,y)\unsure{=}g\handwritten{\phi}(x,y),\quad\handwritten{\phi}(x,gy)\unsure{=}g\handwritten{\phi}(x,y),\quad\handwritten{\phi}(}
\]};
\node[anchor=north west,text width=10.18cm] at (230,660) {\[
\handwritten{\phi}(gx,y)=\handwritten{\phi}(x,gy)=g\handwritten{\phi}(x,y)\qquad(x,y\handwritten{\in}E(S),\ g\handwritten{\in}G(S))
\tag*{(2.1.1.1)}
\]};

\node[anchor=north west,text width=10.18cm,align=justify] at (233,751) {et des relations analogues en $\handwritten{\psi}$, s'expriment alors par l\deleted{a}'associativité et la commutativité de \unsure{\deleted{s}c|\deleted{c}s}es lois de composition partielles, et par la relation};
% The label stands above the equation's baseline; both source curves are drawn, without arrowheads.
\node[anchor=north west,text width=10.18cm] at (235,910) {\[\vphantom{\phi(x)}\tag*{(2.1.1.2)}\]};
\node[anchor=north west,text width=5.1cm] at (405,951) {\[
\handwritten{\phi}(\handwritten{\psi}(x,y),\handwritten{\psi}(z,t))=\handwritten{\psi}(\handwritten{\phi}(x,z),\handwritten{\phi}(y,t))
\]};
\draw[line width=.4pt] (300,1006) .. controls (285,979) and (477,940) .. (563,941) -- (566,929);
\draw[line width=.45pt] (438,934) .. controls (468,958) and (380,977) .. (394,997) .. controls (420,1028) and (427,1023) .. (573,1022) -- (575,1035);
\node[anchor=north west,text width=10.18cm,align=justify] at (235,1038) {pour $x,y,z,t$ respectivement dans $E_{p,q}(S),E_{p,q'}(S),E_{p',q}(S),E_{p',q'}(S)$.};

\node[anchor=north west] at (355,1086) {\inserted{\handwritten{\uline{Sections unité}.}}};
\draw[line width=.4pt] (419,1166) .. controls (409,1151) and (417,1138) .. (434,1135) .. controls (494,1135) and (615,1137) .. (648,1128) .. controls (657,1123) and (658,1116) .. (655,1107);
\node[anchor=north west,text width=10.18cm,align=justify] at (235,1123) {2.2.\hspace{.91cm}Il y a lieu de considérer la section unité de $\mathbb{E}\deleted{t}$ considéré comme extension de $Q_P$ par $G_P$,};
\node[anchor=north west,text width=10.18cm] at (235,1247) {\[
 e_{E/P}\handwritten{\in\Gamma}(E/P)=\mathrm{Hom}_P(P,E)\ ,
\tag*{(2.2.1)}
\]};
\node[anchor=north west,text width=10.18cm,align=justify] at (235,1333) {et de façon symétrique la section unité de $E$, considéré comme extension de $P_Q$ par $G_Q$,};
\node[anchor=north west,text width=10.18cm] at (235,1454) {\[
 e_{E/Q}\handwritten{\in\Gamma}(E/Q)=\mathrm{Hom}_Q(Q,E)\ .
\tag*{(2.2.2)}
\]};

\node[anchor=north west,text width=10.18cm,align=justify] at (235,1543) {En particulier, \uline{le torseur} $E$ \uline{au dessus de} $P\times Q$ (de \unsure{\deleted{g}G|\deleted{G}g}roupe $G_{P\times Q}$) \uline{est canoniquement trivialisé au dessus de} $1_P\times Q$ \uline{par} $e_{E/P}$, \uline{et au dessus de} $P\times1_Q$ \uline{par} $e_{E/Q}$. \uline{\unsure{\deleted{C}c|\deleted{c}C}es deux trivialisations coïncident au dessus de la section unité} $1_P\times1_Q$ \uline{de} $P\times Q$, et définissent une même section};
\node[anchor=north west] at (590,1649) {\inserted{De plus,}};
\node[anchor=north west,font=\fontsize{5}{6}\selectfont] at (591,1681) {\deleted{\handwritten{\illegible{1}}}};
\draw[line width=.5pt] (576,1735) .. controls (566,1703) and (578,1689) .. (600,1688) .. controls (655,1689) and (758,1704) .. (767,1673);
\node[anchor=north west] at (1071,1727) {\inserted{\handwritten{donc}}};
\draw[line width=.4pt] (1086,1782) .. controls (1080,1765) and (1103,1764) .. (1128,1757) .. controls (1146,1750) and (1158,1744) .. (1153,1736);
\node[anchor=north west,text width=10.18cm] at (235,1810) {\[
 e_{E\deleted{/S}}\handwritten{\in\Gamma}(E\handwritten{)}\deleted{\unsure{/S)|/P)|/Q)}}\ ,
\tag*{(2.2.3)}
\]};
\node[anchor=north west,text width=10.18cm,align=justify] at (239,1896) {qu'on appellera par abus de langage la \uline{section unité} (ou \uline{section neutre}) de $E$. Pour prouver l'égalité annoncée\deleted{, et une conséquence} des deux sections $e,e'$ de $E$ déduites respectivement de $e_{E/P}$ et $e_{E/Q}$, on note};

% Long cancellation strokes, as opposed to the independent local corrections above.
\draw[line width=.32pt] (223,653) .. controls (249,925) and (389,1726) .. (434,2142);
\draw[line width=.45pt] (494,189) .. controls (589,294) and (1184,1633) .. (1258,2090);
\draw[line width=.35pt] (939,239) .. controls (1119,601) and (1416,1115) .. (1494,1465);
\draw[line width=.48pt] (1457,352) .. controls (1325,444) and (311,1519) .. (148,1685);
\draw[line width=.4pt] (1555,999) -- (562,2106);
\draw[line width=.4pt] (1560,1425) .. controls (1635,1474) and (1285,1975) .. (1122,2105);
\end{tikzpicture}}
\endgroup
% archive sheet number: 26
